\scanpage{181}
% archive sheet number: 180

Modules discrets \deleted{\unsure{sur}} de type fini sur
\[
A=\deleted{\unsure{W}\times}\lim_{\longleftarrow\,n}\unsure{W|\mathbb W}_{n\,\sigma}[F,\unsure{V}]/(V^n,FV-p)
\]

(Catégorie $C_1$) on considère $C_1^\circ=C_2$.

Et la sous-catégorie \inserted{pleine} $C'_1$ de $C_1^{\unsure{\alpha|\circ}}$ formée des modules locaux annulés par une puissance de $F$, [qui en $C_1^\circ$ correspond \unsure{aux} \deleted{\unsure{limites projectives}} \inserted{à la sous-catégorie $C'_2$} des groupes finis \uline{radiciels} unipotents].

\margin{$C_1$ \deleted{Modules de}\par $C_2$}
\margin{$C'_1$\par $\unsure{\wr|21}$\par $C'_2$}

La dualité de Cartier des groupes finis, définie \deleted{par} sur les modules de type \unsure{cocycliques|cocyclique} définit un \uline{isomorphisme} $C'_1\simeq C'_2$.
D'ailleurs, \unsure{tt objet} $G$ de $C_2$ a un sous-objet canonique tel que $G/G'$ soit le plus grand quotient dans $C'_2$, et de $=$ tout objet de $C_1$ a un plus grand sous-objet \deleted{\unsure{en plus}} dans $C'_1$. Un objet de la nouvelle catégorie $C$ serait un \deleted{\unsure{couple}} \inserted{triple} formé de \deleted{\unsure{un}} $(\deleted{\unsure{Q}}X,Y,u)$, où $X\in C_2$, $Y\in C_1$, et $u$ est un isom de $\alpha(X)$ sur $\beta(Y)$.

\diagram{\begin{tikzcd}[ampersand replacement=\&]
C_1 \arrow[dr,"\alpha"] \& \\
\& C_0 \\
C_2 \arrow[ur,"\beta"'] \&
\end{tikzcd}}
